% Part: many-valued-logic
% Chapter: syntax-and-semantics
% Section: sublogics

\documentclass[../../../include/open-logic-section]{subfiles}

\begin{document}

\olfileid[hi]{mvl}{syn}{sub}

\olsection{~$\LogCL$ के उपतर्कशास्त्रों के रूप में बहुमूल्य तर्कशास्त्र}

सभी प्रचलित बहुमूल्य तर्कशास्त्र ऐसे आव्यूहों से परिभाषित होते हैं जिनमें
$\{\True, \False\}$ के तर्कों पर सत्य-फलन का मान चिरसम्मत सत्य-फलनों से मेल
खाता है। विशेषतः, इन तर्कशास्त्रों में यदि $x \in \{\True, \False\}$, तो
$\tf{\lnot}[\Log L](x) = \tf{\lnot}[\LogCL](x)$; और $\star$, ~$\land$,
$\lor$, $\lif$ में से कोई भी हो, यदि $x, y \in \{\True, \False\}$, तो
$\tf{\star}[\Log L](x,y) = \tf{\star}[\LogCL](x,y)$। दूसरे शब्दों में,
$\{\True,\False\}$ तक प्रतिबंधित $\lnot$, $\land$, $\lor$, $\lif$ के
सत्य-फलन ठीक चिरसम्मत सत्य-फलन हैं।

\begin{prop}\ollabel{prop:mvl-cl}
  मान लें किसी बहुमूल्य तर्कशास्त्र~$\Log L$ की भाषा में संयोजक $\lnot$,
  $\land$, $\lor$, $\lif$ हैं, $\True, \False \in V$, और उसके सत्य-फलन
  निम्न को संतुष्ट करते हैं:
  \begin{enumerate}
    \item\ollabel{prop:not} $\tf{\lnot}[\Log L](x) = \tf{\lnot}[\LogCL](x)$ यदि $x =
    \True$ अथवा $x = \False$;
    \item\ollabel{prop:land} $\tf{\land}[\Log L](x,y) = \tf{\land}[\LogCL](x,y)$,
    \item\ollabel{prop:lor} $\tf{\lor}[\Log L](x,y) = \tf{\lor}[\LogCL](x,y)$,
    \item\ollabel{prop:lif} $\tf{\lif}[\Log L](x,y) = \tf{\lif}[\LogCL](x,y)$,
    यदि $x, y \in \{\True, \False\}$।
  \end{enumerate}
  तब~$V$ में ऐसे किसी भी मूल्यांकन $\pAssign v$ के लिए कि $\pAssign v(p)
  \in \{\True,\False\}$, ~$\pValue v[\Log L](!A) = \pValue
  v[\LogCL](!A)$।
\end{prop}

\begin{proof}
  ~$!A$ पर आगमन द्वारा।
  \begin{enumerate}
  \item यदि $!A \ident p$ परमाण्विक है, तो $\pValue v[\Log L](!A) =
  \pAssign v(p) = \pValue
  v[\LogCL](!A)$.
  \item यदि $!A \ident \lnot B$, तो
  \begin{align*}
    \pValue v[\Log L](!A) & = \tf{\lnot}[\Log L](\pValue v[\Log L](!B)) & &\text{\olref[val]{defn:pValue} से}\\
    & = \tf{\lnot}[\Log L](\pValue v[\LogCL](!B)) && \text{आगमन-परिकल्पना से}\\
    & = \tf{\lnot}[\LogCL](\pValue v[\LogCL](!B)) 
&&\text{परिकल्पना \olref{prop:not} से,}\\
&&&\text{क्योंकि $\pValue v[\LogCL](!B) \in \{\True, \False\}$,}\\
& = \pValue v[\LogCL](!A) &&\text{\olref[val]{defn:pValue} से}।
\end{align*}
\item यदि $!A \ident (!B \land !C)$, तो
\begin{align*}
  \pValue v[\Log L](!A) & = \tf{\land}[\Log L](\pValue v[\Log L](!B), \pValue v[\Log L](!C)) & &\text{\olref[val]{defn:pValue} से}\\
  & = \tf{\land}[\Log L](\pValue v[\LogCL](!B),\pValue v[\LogCL](!C)) && \text{आगमन-परिकल्पना से}\\
  & = \tf{\land}[\LogCL](\pValue v[\LogCL](!B),\pValue v[\LogCL](!C)) 
&&\text{परिकल्पना \olref{prop:land} से,}\\
&&&\text{क्योंकि $\pValue v[\LogCL](!B),\pValue v[\LogCL](!C) \in \{\True, \False\}$,}\\
& = \pValue v[\LogCL](!A) &&\text{\olref[val]{defn:pValue} से}।
\end{align*}
\end{enumerate}
$!A \ident (!B \lor !C)$ और $!A \ident (!B \lif !C)$ वाली स्थितियाँ
समान हैं।
\end{proof}

\begin{cor}
  यदि कोई बहुमूल्य तर्कशास्त्र \olref{prop:mvl-cl} की शर्तें संतुष्ट करता
  है, $\True \in V^+$ और $\False \notin V^+$, तो
  ${\Entails[\Log L]} \subseteq {\Entails[\LogCL]}$, अर्थात् यदि $\Gamma
  \Entails[\Log L] !B$, तो $\Gamma \Entails[\LogCL] !B$। विशेषतः,
  $\Log L$ की प्रत्येक सर्वसत्यता चिरसम्मत सर्वसत्यता भी है।
\end{cor}

\begin{proof}
  हम प्रतिधनात्मक सिद्ध करते हैं। मान लें $\Gamma \Entails/[\LogCL] !B$।
  तब कोई ऐसा !!{valuation}~$\pAssign v\colon \PVar \to \{\True,
  \False\}$ है कि सभी $!A \in \Gamma$ के लिए $\pValue v[\LogCL](!A) =
  \True$ और $\pValue v[\LogCL](!B) = \False$। चूँकि $\True, \False \in
  V$, इसलिए !!{valuation}~$\pAssign v$, ~$\Log L$ के लिए भी
  !!a{valuation} है। \olref{prop:mvl-cl} से सभी $!A \in \Gamma$ के लिए
  $\pValue v[\Log L](!A) = \True$ और $\pValue v[\Log L](!B) = \False$।
  चूँकि $\True \in V^+$ और $\False \notin V^+$, इसका अर्थ है $\pAssign v
  \Entails[\Log L] \Gamma$ और $\pAssign v \Entails/[\Log L] !B$, अर्थात्
  $\Gamma \Entails/[\Log L] !B$।
\end{proof}
\end{document}
